\clearpage
\SourceLeaf{147}
\begin{center}
\textsc{Gauss} an \textsc{Olbers}. Göttingen, 24. Januar 1812.
\end{center}

\noindent . . . . . . . . . . . Herr\ \textsc{Delambre} soll, wie ich höre, im französischen\ Moniteur in seiner Art eine sehr weitschweifige Chrie über die moindres carrés gegeben haben. Ich lese den Moniteur nicht, der auch nur in jährlichen Lieferungen hierher kommt. Vielleicht finden Sie einmal Gelegenheit, \textls{öffentlich} zu bezeugen, dass ich Ihnen schon bei unserer ersten persönlichen Begegnung im Jahr 1803 die Hauptpunkte davon dargelegt habe. Unter meinen Papieren finde ich, dass ich im Juni 1798, als ich jene Methode bereits \textls{lange} angewandt hatte, zuerst \textsc{Laplaces} Methode gesehen und ihre Unvereinbarkeit mit den Grundsätzen der Wahrscheinlichkeitsrechnung in einem kurzen Notizjournal\HistoricalNote{[*)]}{[*)}{In dem bereits auf Seite 20 erwähnten Tagebuch von \textsc{Gauss} findet sich aus dem Jahr 1798 die Notiz: »\emph{Calculus probabilitatis contra La Place defensus. Gott. Jun. 17}«.]} über meine mathematischen Beschäftigungen festgehalten habe. Im Herbst 1802 habe ich die \textsc{viii}\textsuperscript{ten} Ceres-Elemente in meinem astronomischen Entwurfsbuch eingetragen, nachdem ich sie nach der Methode der kleinsten Quadrate gefunden hatte. Die Papiere, in denen ich in früheren Jahren, z.\,B.\ im Frühjahr 1799, jene Methode auf \textsc{Ulugh Beighs} Zeitgleichungstafel angewandt habe, sind verloren gegangen. Das Einzige, worüber man sich wundern kann, ist, dass dieses Prinzip, das sich so leicht von selbst darbietet, dass man auf den Gedanken allein gar keinen besonderen Wert legen kann, nicht schon 50 oder 100 Jahre früher von anderen, z.\,B.\ von \textsc{Euler} oder \textsc{Lambert} oder \textsc{Halley} oder \textsc{Tobias Mayer}, angewandt worden ist. Doch es kann ja sehr wohl sein, dass z.\,B.\ letzterer so etwas angewandt hat, ohne es zu verkünden, so wie jeder Rechner sich notwendigerweise selbst eine Menge Vorteile und Methoden verschafft, die er nur gelegentlich mündlich weitergibt.

\noindent . . . . . . . . . . . . .

